% Moved 2026-10-28 from dl_rnn.tex: how to get one vector from a sequence, with any recurrent cell.
% Follows the plain RNN, GRU and LSTM sections (dl_rnn, dl_gru, dl_lstm).
%%%%%%%%%%%%%%%%%%%%%%%%%%%%%%%%%%%%%%%%%%%%%%%%%%%
\begin{frame}[fragile] \frametitle{One Vector per Sequence: Three Readouts}
\begin{itemize}
% [upgrade-deck: rewritten 2026-10-26 | Task 5c | long bullets split into a lead plus sub-points; every fact kept; original below]
% \item For a sentiment label or a fault class we need one vector for the whole sequence, then a classifier on top. Three common ways to get it:
% \item \textbf{Last state} $h_T$: simple, but the whole sequence is squeezed through one vector. Order is not lost, but old information has faded (recency bias)
% \item \textbf{Pooling over time}: keep every $h_t$ and take the mean or the maximum of each feature over all steps. An early, strong signal survives
% \item \textbf{Attention}: keep every $h_t$ and let the model learn how much each step counts for this decision
% \item A fourth, orthogonal fix is to read the sequence in both directions (next slides)
\item With a plain RNN, a GRU or an LSTM, a sentiment label or a fault class needs one vector for the whole sequence, then a classifier on top. Three common ways:
\item \textbf{Last state} $h_T$: simple, but the whole sequence is squeezed through one vector
  \begin{itemize}
  \item Order is not lost, but old information has faded (recency bias: the newest steps dominate)
  \end{itemize}
\item \textbf{Pooling over time}: take the mean or the maximum of each feature over all the $h_t$. An early, strong signal survives
\item \textbf{Attention}: let the model learn how much each $h_t$ counts for this decision
\item A fourth, separate fix: read the sequence in both directions (the Bidirectional RNN frame)
\end{itemize}
\begin{center}
\adjustbox{max width=\linewidth, max totalheight=0.28\textheight}{%
\begin{tikzpicture}[>=latex, every node/.style={font=\scriptsize}, st/.style={draw, circle, minimum size=4mm, inner sep=0pt}, bx/.style={draw, rounded corners, fill=lightnyupurple!30, minimum width=19mm, minimum height=6mm}]
\begin{scope}
\foreach \i/\x in {1/0, 2/0.7, 3/1.4, 4/2.1} \node[st] (a\i) at (\x,0) {};
\node[anchor=south] at (a1.north) {$h_1$};
\node[anchor=south] at (a4.north) {$h_T$};
\node[bx] (ab) at (1.05,-1.3) {last state};
\draw[->, thick] (a4) -- (ab);
\end{scope}
\begin{scope}[xshift=3.6cm]
\foreach \i/\x in {1/0, 2/0.7, 3/1.4, 4/2.1} \node[st] (b\i) at (\x,0) {};
\node[anchor=south] at (b1.north) {$h_1$};
\node[anchor=south] at (b4.north) {$h_T$};
\node[bx] (bb) at (1.05,-1.3) {mean or max};
\foreach \i in {1,2,3,4} \draw[->] (b\i) -- (bb);
\end{scope}
\begin{scope}[xshift=7.2cm]
\foreach \i/\x in {1/0, 2/0.7, 3/1.4, 4/2.1} \node[st] (c\i) at (\x,0) {};
\node[anchor=south] at (c1.north) {$h_1$};
\node[anchor=south] at (c4.north) {$h_T$};
\node[bx] (cb) at (1.05,-1.3) {weighted sum};
\draw[->, line width=0.4pt] (c1) -- (cb);
\draw[->, line width=1.6pt] (c2) -- (cb);
\draw[->, line width=0.4pt] (c3) -- (cb);
\draw[->, line width=1.0pt] (c4) -- (cb);
\end{scope}
\end{tikzpicture}}

{\tiny Thicker arrow, larger share. Last state uses one step; pooling treats all steps alike; attention learns the shares.}
\end{center}
\end{frame}

%%%%%%%%%%%%%%%%%%%%%%%%%%%%%%%%%%%%%%%%%%%%%%%%%%%
\begin{frame}[fragile] \frametitle{Early Evidence: Which Readout Works?}
\begin{itemize}
\item Task: sequences of 60 steps of noise; the label is whether a $+3$ marker sits in steps 1 to 3. Chance is 50\%. One run each
\end{itemize}
\begin{center}
\small
\adjustbox{max width=\linewidth}{%
\begin{tabular}{lrc}
\hline
Model & Parameters & Test accuracy \\
\hline
Plain RNN, last state & 338 & 85.3\% \\
Plain RNN, mean over time & 338 & 90.4\% \\
Plain RNN, max over time & 338 & 99.2\% \\
Plain RNN, bidirectional & 674 & 99.1\% \\
GRU, max over time & 946 & 99.3\% \\
GRU, bidirectional & 1,890 & 99.2\% \\
GRU, last state (100 epochs) & 946 & 52.2\% \\
LSTM, last state (100 epochs) & 1,250 & 50.9\% \\
\hline
\end{tabular}}
\end{center}
\begin{itemize}
% [upgrade-deck: rewritten 2026-10-26 | Task 5c | long bullets split into a lead plus sub-points; every fact kept; original below]
% \item Max pooling brings the early marker straight to the classifier. Bidirectional works because the backward pass ends at step 1, where the marker is fresh
% \item The last-state GRU and LSTM did not learn this task in our training budget, while the plain RNN did partly. We do not claim a ranking of cell types from one synthetic run
\item \textbf{Max pooling:} brings the early marker straight to the classifier
\item \textbf{Bidirectional:} the backward pass ends at step 1, where the marker is fresh
\item The last-state GRU and LSTM did not learn this task in our training budget, while the plain RNN did partly
\item We do not claim a ranking of cell types from one synthetic run
\end{itemize}
\end{frame}

%%%%%%%%%%%%%%%%%%%%%%%%%%%%%%%%%%%%%%%%%%%%%%%%%%%
\begin{frame}[fragile] \frametitle{Bidirectional RNN}
\begin{center}
\adjustbox{max width=0.7\linewidth, max totalheight=0.40\textheight}{\includegraphics{rnnbi}}
\end{center}
\begin{itemize}
\item Two recurrent layers: one reads left to right, one right to left. At each step their states are joined, so every step has context from both sides
\item The cost is twice the parameters (338 became 674 in the test before) and the whole sequence must be available
\item Fine for classifying a finished recording or sentence. Not possible for real-time forecasting, where the future has not happened yet
\end{itemize}
\end{frame}

%%%%%%%%%%%%%%%%%%%%%%%%%%%%%%%%%%%%%%%%%%%%%%%%%%%
% [upgrade-deck: rewritten 2026-10-26 | Tasks 5 and 5c | the definitions got their own frame with an Intuition callout; the worked example follows on the next frame; original below]
% \begin{frame}[fragile] \frametitle{Attention: A Weighted Look Back}
% \begin{itemize}
% \item Keep all hidden states $h_1, \dots, h_T$. Score each against a query $q$, turn scores into weights with softmax, and mix: $s_i = q \cdot h_i$, \quad $\alpha = \mathrm{softmax}(s)$, \quad $c = \sum_i \alpha_i h_i$
% \item Example: $h_1 = (1,0)$, $h_2 = (0,1)$, $h_3 = (1,1)$ and $q = (1,0)$
% \end{itemize}
\begin{frame}[fragile] \frametitle{Attention: A Weighted Look Back}
\begin{itemize}
\item Keep all hidden states $h_1, \dots, h_T$, not only the last one
\item \textbf{Score} each state against a query $q$ (a vector for what we are looking for): $s_i = q \cdot h_i$
\item \textbf{Weights:} softmax turns the scores into positive numbers that add up to 1: $\alpha = \mathrm{softmax}(s)$
\item \textbf{Mix} the states with those weights: $c = \sum_i \alpha_i h_i$
\end{itemize}
\begin{block}{Intuition}
Think of reading your notes with a highlighter. The query is the question you ask, the scores say which lines match it, and the weights decide how much of each line goes into your answer.
\end{block}
\end{frame}

%%%%%%%%%%%%%%%%%%%%%%%%%%%%%%%%%%%%%%%%%%%%%%%%%%%
\begin{frame}[fragile] \frametitle{Attention: A Worked Example}
\begin{itemize}
\item Example: $h_1 = (1,0)$, $h_2 = (0,1)$, $h_3 = (1,1)$ and $q = (1,0)$
\end{itemize}
\begin{center}
\small
\adjustbox{max width=\linewidth}{%
\begin{tabular}{lccc}
\hline
 & $h_1$ & $h_2$ & $h_3$ \\
\hline
score $q \cdot h_i$ & 1 & 0 & 1 \\
weight $\alpha_i$ & 0.4223 & 0.1554 & 0.4223 \\
\hline
\end{tabular}}
\end{center}
\begin{itemize}
\item Context $c = 0.4223\,(1,0) + 0.1554\,(0,1) + 0.4223\,(1,1) = (0.8446,\ 0.5777)$
\item A plain average would give $(0.6667,\ 0.6667)$: attention leans toward the states that match the query
\end{itemize}
\end{frame}

%%%%%%%%%%%%%%%%%%%%%%%%%%%%%%%%%%%%%%%%%%%%%%%%%%%
\begin{frame}[fragile] \frametitle{Attention: The Road to Transformers}
\begin{itemize}
\item Attention was first added on top of RNNs to ease the one-vector bottleneck: the decision can look back at every step, not only the last one
\item A Transformer keeps attention and drops the recurrence. It adds position information so that order is not lost, and it reads all steps at once, which suits parallel hardware
\item That is a pointer, not part of this session: the point here is why the bottleneck mattered
\item Recurrent nets have not disappeared. Small, cheap models that process a stream one step at a time are still a good fit when memory and compute are tight
\end{itemize}
\end{frame}

%%%%%%%%%%%%%%%%%%%%%%%%%%%%%%%%%%%%%%%%%%%%%%%%%%%
\begin{frame}[fragile] \frametitle{Quick Check: Sequences}
\begin{block}{Think About It}
1. A sensor gives 8 readings per step, and a layer has 32 hidden units. The window grows from 100 to 1000 steps. What happens to the parameter count?
\begin{enumerate}
  \item[A.] Both a dense layer on the window and a recurrent layer grow about ten times.
  \item[B.] The dense layer grows about ten times; the recurrent layer stays the same.
  \item[C.] The recurrent layer grows about ten times; the dense layer stays the same.
  \item[D.] Both stay the same.
\end{enumerate}
2. In the early-marker task, which readout lifted a plain RNN from 85\% to about 99\%?
\begin{enumerate}
  \item[A.] Last state only.
  \item[B.] Max over time (or reading in both directions).
  \item[C.] More hidden units with the same last-state readout.
  \item[D.] Shuffling the steps.
\end{enumerate}
\end{block}
\end{frame}

%%%%%%%%%%%%%%%%%%%%%%%%%%%%%%%%%%%%%%%%%%%%%%%%%%%
\begin{frame}[fragile] \frametitle{Quick Check: Sequences}
\begin{block}{Answer}
\textbf{1. Correct: B.} The dense layer on the window has $T \times 8 \times 32$ weights plus biases: 25,632 at 100 steps and 256,032 at 1000. The recurrent layer reuses the same weights at every step: 1,344 for a plain RNN in both cases.

\textbf{2. Correct: B.} Max over time keeps the early marker's strong response, and a backward pass ends exactly at the early steps. We did not test C, and D destroys the order the label depends on.
\end{block}
\end{frame}
